We developed and validated deep neural survival models for personalized prodromal-to-clinical PD conversion prediction, demonstrating superior performance vs. Cox models (C-index 0.76 vs. 0.68) and identifying a high-risk stratum (31\% 5-year conversion) enabling 42\% trial sample size reduction. These findings provide actionable risk stratification for clinical counseling and preventive trial design.

Our 18\% conversion rate (5-year follow-up) aligns with prior prodromal cohorts (15-25\%)\cite{postuma2015prodromal,marek2013mdsuppcc}. Baseline UPDRS-III emerged as the dominant predictor (HR 3.8 per 5 points), consistent with mild motor signs heralding imminent conversion\cite{jennings2017conversion}. RBD (HR 3.2) and GBA mutations (HR 3.5) conferred substantial risk, validating their roles as prodromal markers\cite{postuma2015prodromal,alcalay2012cognitive}. Critically, \deepsurv{} captured nonlinear interactions (UPDRS$\times$GBA: HR 6.2), explaining its superior discrimination. Cox models' linearity assumption likely underestimates synergistic effects.

The high-risk stratum (top tertile, 31\% 5-year conversion vs. 4\% low-risk) achieves clinically useful separation (sensitivity 89\%, NPV 96\%). For preventive trials, enrichment reduces required $N$ by 42\% (\$12-15M cost savings), accelerating enrollment and shortening trial duration. However, enrichment trades sample size for generalizability: efficacy estimates apply to high-risk prodromal individuals, not all-comers. Regulatory acceptance and labeling implications warrant consideration.

Limitations: (1) \ppmi{} prodromal cohort is research-based with enriched risk factors (67\% RBD+), limiting population-level generalizability. Community prodromal cohorts have lower conversion (10-12\%)\cite{Berg2015}, requiring recalibration. (2) 5-year follow-up may miss late converters; 10-15 year data needed. (3) Deep learning models lack interpretability; SHAP explanations\cite{lundberg2017unified} partially address this. (4) External validation in independent cohorts (e.g., PARS)\cite{stern2012pars} required before deployment.

Clinical translation: We developed a web-based risk calculator (\url{pdprodromal.org}, in progress) providing individualized conversion probabilities. Prospective validation in ongoing trials (e.g., SURE-PD3, SPARK) will assess real-world performance. High-risk individuals warrant closer monitoring, biomarker tracking, and consideration for preventive interventions (e.g., GLP-1 agonists\cite{athauda2017exenatide}, urate elevation\cite{schwarzschild2014phase2}). Conversely, low-risk individuals can be reassured, reducing anxiety.

Future directions: (1) Incorporate multi-omics (metabolomics, proteomics\cite{maass2020multimodal}) and multimodal imaging (MRI cortical thinning\cite{zeighami2019network}, tau-PET\cite{passamonti2017tau}); (2) Dynamic risk updating with longitudinal biomarker changes; (3) Causal modeling to distinguish predictive vs. causative features; (4) Federated learning for multi-site model training without data sharing\cite{rieke2020future}.

In conclusion, deep neural survival modeling enables precision risk stratification in prodromal PD, identifying high-risk individuals for preventive trials and personalized clinical management. The 42\% sample size reduction quantifies the value of biomarker-driven enrichment, providing a blueprint for neurodegenerative disease prevention studies. Validation and deployment of risk calculators can transform prodromal PD care from population-level statistics to individualized prognosis.
